\usepackage{pgfplots}
\usepackage{pgfplotstable}
%% ==================== DOCUMENT COLORS ====================
% Define global colors for hyperlinks and URLs
\definecolor{linkborderColor}{rgb}{0.1329,0.247,0.467}   % Blue for link borders
\definecolor{urlborderColor}{rgb}{0,1,1}                % Cyan for URL borders
\definecolor{urlColor}{rgb}{1,0,0}                      % Red for URL text

%% ==================== HYPERREF SETUP ====================
\hypersetup{
%author = {Thomas AUBIN},                            % Uncomment to set PDF author
%pdftitle = {Nom du Projet - Sujet du Projet},       % Uncomment to set PDF title
    pdfsubject = {Mémoire de Projet},                    % PDF document subject
    pdfkeywords = {Tag1, Tag2, Tag3, \dots},             % PDF keywords for indexing
    pdfstartview={FitH},                                 % Default PDF view (fit width)
    linkbordercolor = {linkborderColor}                 % Color for hyperlink borders
%urlbordercolor = \color{urlborderColor},           % Uncomment to set URL border color
%urlcolor = \color{urlColor}                        % Uncomment to set URL text color
}


%% Divergence palette for heatmaps

\definecolor{errmax}{RGB}{215,48,39}   % rouge
\definecolor{neu}{RGB}{255,255,255} % blanc
\definecolor{errmin}{RGB}{69,117,180}  % bleu

\pgfplotsset{compat=1.18}

% Commandes pour définir les seuils min, milieu et max de la coloration
\newcommand{\setdivergentmin}[1]{\def\divergentminval{#1}}
\newcommand{\setdivergentmid}[1]{\def\divergentmidval{#1}}
\newcommand{\setdivergentmax}[1]{\def\divergentmaxval{#1}}

% Définir des valeurs par défaut
\setdivergentmin{0}
\setdivergentmid{3}
\setdivergentmax{12}

\newcommand{\divergentcell}[1]{%
    \pgfmathparse{
        #1 < \divergentmidval ?
        (min(100, max(0, (\divergentmidval - #1) / (\divergentmidval - \divergentminval) * 100))) :
        (min(100, max(0, (#1 - \divergentmidval) / (\divergentmaxval - \divergentmidval) * 100)))
    }%
    \let\percent=\pgfmathresult
    \ifdim #1pt < \divergentmidval pt
        \edef\temp{\noexpand\cellcolor{errmin!\percent!neu}}%
    \else
        \edef\temp{\noexpand\cellcolor{errmax!\percent!neu}}%
    \fi
    \temp
    #1%
}

% Commandes pour la coloration à 2 niveaux (ex: blanc à rouge)
\newcommand{\settwotonemin}[1]{\def\twotoneminval{#1}}
\newcommand{\settwotonemax}[1]{\def\twotonemaxval{#1}}

% Définir des valeurs par défaut
\settwotonemin{0}
\settwotonemax{5}

\newcommand{\twotonecell}[1]{%
    % Calculer le pourcentage de la valeur dans l'intervalle [min, max]
    \pgfmathparse{min(100, max(0, ((#1 - \twotoneminval) / (\twotonemaxval - \twotoneminval)) * 100))}%
    \let\percent=\pgfmathresult
    % Appliquer la couleur de neu (blanc) à errmax (rouge)
    \edef\temp{\noexpand\cellcolor{errmax!\percent!neu}}%
    \temp
    #1%
}
